\subsubsection{Nested Loops Join (NLJ)}
\begin{algorithm}
    \caption{Nested Loops Join Algorithm}
    \begin{algorithmic}[1]
        \Procedure{NestedLoopsJoin}{$R$, $S$}
            \For{each page $p_R \in R$}
                \For{each page $p_S \in S$}
                    \State Join the tuples on page $p_R$ with the tuples on page $p_S$ (evaluate the predicate for all pairs)
                    \State Add matching tuples to the join output
                \EndFor
            \EndFor
        \EndProcedure
    \end{algorithmic}
\end{algorithm}

The cost here is \cost{$P_R + P_R \cdot P_S$ I/Os}, where $P_X$ is the number of pages in the relation $X$.
Furthermore, the DBMS should put the smaller relation on the outer loop to improve performance.

For multiple joins, we want to start with the smallest table (on the inside) and move up to larger ones, processing the largest one \textit{last} (such that it is the outer loop)
